\newpage
\section{Resumo}

Este trabalho apresenta a simulação de um manipulador industrial Stäubli TX90, de
seis graus de liberdade, executando a pintura da bandeira do Brasil, como aplicação
integrada das técnicas de robótica estudadas na disciplina. A partir da definição
geométrica das formas da bandeira, gerou-se a trajetória cartesiana do efetuador com
perfil de velocidade LSPB --- no modo de traçado dos contornos ou de preenchimento
por varredura em serpentina ---, intercalada com deslocamentos livres a uma estação
de troca de cores. Cada ponto da trajetória foi convertido para o espaço articular
por cinemática inversa, validada por cinemática direta e pela análise do menor valor
singular do Jacobiano, que mede a proximidade de singularidades. Em seguida, a
dinâmica inversa forneceu os torques exigidos, confrontados com os limites do
manipulador, e simularam-se dois controladores em malha fechada: PD com compensação
de gravidade e torque computado. Os resultados confirmam a viabilidade da tarefa:
todas as poses admitem solução de cinemática inversa longe de singularidades, os
limites de velocidade e de torque são respeitados e ambos os controladores
acompanham a referência com erros reduzidos, com desempenho superior do torque
computado. Por fim, discutem-se as limitações do modelo e os passos necessários para
uma eventual implementação física.